%--------------------------------------------------------
\section{Organisatorisches}

\subsection{Was bedeutet eigenständiges Arbeiten? (Vgl. \cite{Menge})} 

\paragraph{Kommunikation}
Die Eigenständigkeit ist nicht gleichzusetzen mit fehlender Kommunikation. Es ist wichtig, dass du dich regelmäßig mit deinen Gutachter\_innen in Verbindung setzt und über den
Fortgang deiner Arbeit berichtest. Dies hat die folgenden Gründe:
\begin{itemize}
    \item Es lassen sich Fragen in Bezug auf die Aufgabenstellung oder in Bezug auf mögliche Lösungen klären.
    \item Nicht selten stockt eine Abschlussarbeit, weil ein Problem besteht, das sich in einem Gespräch schnell lösen ließe. Solltest du nicht weiterkommen, nimm deshalb schnell Kontakt mit deiner Betreuerin/ deinem Betreuer auf. Fragen zu stellen und sich von anderen inspirieren zu lassen gehört zum Alltag von Ingenieur\_innens und wird nicht als Makel angesehen.
    \item Durch regelmäßige Kommunikation wird sichergestellt, dass du das Thema nicht verfehlst.
\end{itemize}

\paragraph{Strukturiert arbeiten}
\label{sec:struktur}

Gleich zu Beginn deiner Arbeit solltest du eine grobe Zeitplanung, z.B. in Form eines Balkenplanes und eines Meilensteinplans, anlegen und diese während der Arbeit ständig korrigieren und fortschreiben. Durch solche einfachen Mittel des Projektmanagements kannst du unnötigen Zeitverlust vermeiden, der oft dadurch entsteht, dass man sich mit Randproblemen zu lange aufhält. Spreche diese Zeitpläne mit deiner Betreuerin/ deinem Betreuer durch.

Dokumentiere deinen Fortschritt. Durch regelmäßige Berichte dokumentierst du, dass Sie das Thema bearbeitest und auf welche Weise du das tust. Dies ist für die Bewertung des praktischen Teils von großer Bedeutung.

%Sie sollten  regelmäßig den aktuellen Stand der Arbeit präsentieren können. Auch sollten Sie den Fortschritt aufzeigen können, der sich zwischen dem vorangehenden und dem aktuellen Termin ergeben hat. Als Diskussionsgrundlage ist es sinnvoll, Bilder, Skizzen, Prototypen, Listings usw. zu den Treffen mitzubringen.


\subsection{Was ist ein Exposé?}

Zur Unterstützung bei der Themenfindung sowie des Arbeits- und Schreibprozesses empfehle ich ein Exposé zu erstellen. Es wird  zu Beginn der Bearbeitung erstellt und klärt Schwerpunkte und zeitliche Eckpunkte der Abschlussarbeit. Im Exposé wird die Arbeit präzisiert, nachdem die ersten Recherchen und Strukturierungen vorgenommen wurden. 

Das Exposé enthält folgende Teile:
\begin{itemize}
\item Thema (Titel, Untertitel, Abgrenzung),
\item Motivation,
\item Forschungsstand (bisherige Recherche, vorläufiges Fazit),
\item Zentrale Fragestellung,
\item Methodisches Vorgehen (Literaturrecherche, Berechnungen, Auswertungen etc.),
\item Gegenwärtige Gliederung.
\end{itemize}
Das Exposé sollte ca. $1$ bis $3$ Seiten lang sein.

\subsection{Literaturrecherche}

Zum Stand der Technik suchst du am besten Informationen und Literatur in der Bibliothek oder im Internet. Hier empfehle ich folgende Portale:
\begin{itemize}
\item die Ratschläge such Literaturrecherche der Bibliothek der HTW\\ 
\hyperlink{https://bibliothek.htw-berlin.de/literatur-suchen/}{\emph{HTW Bibliothek online}}
    \item \hyperlink{scholar.google.com/}{\emph{Google Scholar}}
    \item  
    \hyperlink{http://ezb.uni-regensburg.de/fl.phtml?bibid=AAAAA&colors=7&lang=de&notation=ZN}{\emph{Elektronische Zeitschriftenbibliothek (EZB) der HTW}}
    \item     insbesonder die IEEE-Datenbank unter 
    \hyperlink{https://ieeexplore.ieee.org/xpl/RecentIssue.jsp?punumber=2}{\emph{IEEE Explore}}
\end{itemize}


\subsection{Versionierung}



Mit der Nutzung eines Versionierungssystems kannst du
\begin{itemize}
    \item Änderungen an einer Datei oder einer Gruppe von Dateien im Laufe der Zeit aufzeichnen
    \item ausgewählte Dateien auf einen früheren Zustand zurücksetzen
    \item das gesamte Projekt in einen früheren Zustand zurückversetzen 
    \item Änderungen im Zeitverlauf vergleichen
    \item 
    
\end{itemize}
% Ich empfehle GitHub. \hyperref{https://git-scm.com/book/en/v2/Getting-Started-About-Version-Contro}{cc}
 
\begin{itemize}

\item rFree GitHub/ GitLab/ Bitbucket: 
Suitable for open-source projects. Not suitable for sensitive files 
Public or private repositories.

\item rhttps://github.com/
(mit registierung] in overleaf github synchronisierung aktivieren

\item reinmal pushen lokal clone respository new rository für ssimulastins/ xperimenuierdaten

\item regelmäßig pushen und lokal pullen mit aussagekräftigen textgen

\item rim repsitory kann man hiostory und Änderungen anschaen und auf alte cersionen wechseln
    
\end{itemize}

\subsection[Zuverlässige Ergebnisse erzeugen]{Zuverlässige Ergebnisse erzeugen: Wie du Entwürfe, Gleichungen oder Code kritisch überprüfst}

Vorsicht vor zu optimistischerVoreingenoommenheit


\begin{enumerate}
    \item Wann immer Sie eine Lösung gefunden haben, sollten Sie ihr nicht vertrauen. Stellen Sie sie stattdessen in Frage und versuchen Sie, sie zu widerlegen. Führen Sie Plausibilitätsprüfungen durch. 
\item Obwohl es sehr hilfreich ist, wenn andere (Kollegen, Vorgesetzte) Ihre Arbeit gegenprüfen, verlassen Sie sich nicht darauf, dass andere alle Ihre Fehler entdecken. Eine Gruppe kann sich leicht auf einen falschen Ansatz kommen und so ein falsches Gefühl von Sicherheit und gemeinsamer Verantwortung vermitteln. Deshalb ist
Feedback ist wichtig, aber nicht ausreichend.
\item Schreiben Sie keine nett klingenden, aber vagen Worte, die so klingen sollen, als ob Sie verstehen, was wenn Sie in Wirklichkeit nicht so sicher sind. Machen Sie stattdessen deutlich, wie groß Ihr Vertrauen in die erzielten Ergebnisse ist. Seien Sie einfach ehrlich. Wenn Sie ein Konzept nicht verstanden haben und keine Ahnung haben, ob die Ergebnisse richtig sind, sagen Sie es einfach. 
\item  Nehmen Sie keine Abkürzungen; stellen Sie sicher, dass das, was Sie liefern, korrekt ist, auch wenn es unvollständig ist.  Zusammenfassend: Seien Sie stets kritisch gegenüber Ihrer eigenen Arbeit und zeigen Sie transparent, wie viel 
Vertrauen in Ihre Ergebnisse.
\end{enumerate}


HOW TO CHECK OBTAINED SOLUTIONS?
\begin{itemize}
    \item Prüfen Sie, ob alle den Formeln zugrunde liegenden Annahmen 
Sie anwenden, erfüllt sind, um sicherzustellen, dass kein 
Fehler im allgemeinen Ansatz.
\item Führen Sie mehrere Plausibilitätsprüfungen durch, um sicherzustellen, dass 
keinen Fehler in der Berechnung gab.
\end{itemize}


SOME GENERIC PLAUSIBILITY CHECKS
\begin{itemize}
\item Einheiten als Schnellprüfung verwenden: Wenn die Einheiten auf beiden Seiten eines Gleichheitszeichens nicht übereinstimmen, ist die Gleichung 
falsch.
\item Verwenden Sie klar getrennte Schreibweisen für Vektoren, um eine schnelle Überprüfung der Korrektheit zu ermöglichen: 
    \item die gleiche Frage mit einem anderen Koordinatensystem beantworten.
    \item  Prüfung von Sonderfällen, in denen die Antwort leichter zu finden ist, z. B. wenn einige Werte oder Ableitungen 0 sind. 
oder Ableitungen 0 sind. Ein intuitiver Spezialfall für dynamische Systeme ist oft das statische Gleichgewicht. 
    \item VorZeichen überprüfen. Wenn eine Variable ansteigt, in welche Richtung wird sich dann eine andere Variable ändern wird?
    \item  einen zuvor integrierten Term erneut differenzieren (um zu überprüfen, ob die Integration korrekt durchgeführt wurde). 
        \item Dimensionen einer Variablen prüfen, z. B. ob Vektoren und Skalare addiert werden?
    \item  verschiedene, alternative Lösungsansätze zu prüfen.
\end{itemize}

SOME FREQUENT MISTAKES IN CALCULATIONS (GENERIC)
\begin{itemize}
    \item Was also, wenn Ihre Plausibilitätsprüfungen fehlschlagen? Einige Fehler treten bei Studenten besonders häufig auf, so ist es wahrscheinlich, dass eine Berechnung bei einer Plausibilitätsprüfung durchgefallen ist, weil: 
    \item Das Koordinatensystem oder die Bezugspunkte waren nicht klar definiert,
    \item formeln wurden angewendet, ohne die diesen Formeln zugrunde liegenden Annahmen zu überprüfen. Der erste Schritt besteht darin, sich bewusst zu machen, dass viele Formeln nur für Spezialfälle gelten, und der zweite Schritt besteht darin gründlich zu prüfen, ob das gegebene Problem wirklich in diesen Spezialfall fällt.
    \item zu wenig Zeit in die Bestimmung der Vorzeichen der Variablen investiert wurde, in der Annahme, dass Vorzeichenfehler 
"kleine Fehler" sind, oder dass sie durch Versuch und Irrtum behoben werden können. Nichts von alledem ist wahr. 
    \item Die Variablen wurden implizit (und inkonsistent) definiert. Variablendefinitionen müssen explizit sein. 
    \item vektoren und Skalare wurden nicht sauber getrennt. 
    \item Einheiten wurden vernachlässigt. 
        \item Numerische Ergebnisse werden ohne Berücksichtigung numerischer Ungenauigkeiten in den Simulationen interpretiert. Wenn 
Wenn möglich, führen Sie auch einige analytische Berechnungen durch. Wenn zum Beispiel eine Zahl in einer Simulation nahe an 
Null erscheint, prüfen Sie, ob es Gründe für die Annahme gibt, dass sie genau Null sein könnte...
\end{itemize}


HOW TO CHECK PLAUSIBILITY
\begin{enumerate}
    \item Stellen Sie eine Hypothese auf, welche Bedingung zutreffen sollte. Schreiben Sie sie auf. 
    \item Prüfen Sie dann, ob die Bedingung tatsächlich erfüllt ist. 
    \item Lassen Sie Schritt 1 NICHT aus. Der Optimismus kann Sie dazu verleiten, zu glauben dass das, was Sie finden, das ist, was Sie die ganze Zeit erwartet haben.
\end{enumerate}
\begin{itemize}
    \item Static equilibrium: Set all derivatives to zero 
        \item Choose special variables/parameters (e.g. zero or infinity)

    \item !check signs: What is the influence of an increase in one variable? In which direction 
does the result change?
    \item !Alternative approach
    \item Derive EoM using Newton-Euler AND using Lagrange
    \item Use another coordinate system in calculations
    \item Inverse calculation 
        \item Integration versus differentiation
    \item Forward versus inverse dynamics
    \item !Dimensions
    \item Check that the product of vectors, matrices, scalars can be calculated
    \item Check you are never dividing by a vector or matrix
Vector Directions
    \item Check that the direction of a vector complies with your expectation, … 
\end{itemize}

HOW TO CREATE A RELIABLE SIMULATION
\begin{enumerate}

    \item 1. Schreiben Sie alle Gleichungen auf Papier aus, überprüfen Sie zunächst die Annahmen und führen Sie Plausibilitätsprüfungen 
manuell durchführen (siehe oben). 
    \item Wandeln Sie Ihre Papierableitungen in Code um. 
    \item Wählen Sie die gleichen Variablennamen in Ihrem Code,
    \item 
    einheitliche Konventionen für die Umwandlung von Symbolen, tiefgestellten und hochgestellten Zeichen usw. in Code verwenden. 
    \item Haben Sie keine Angst vor langen Namen in Ihrem Code (z. B. "gamma" statt "g"). Das verhindert Fehler und macht den Code besser lesbar. 
    \item Make sure you know how to pronounce Greek letters and write them out in code. 
    \item Ableitung und Umsetzung von Bewegungsgleichungen, die mit verschiedenen Methoden (z. B. Lagrange 
und Newton-Euler). Ergebnisse überlagern. 
    \item Führen Sie Plausibilitätsprüfungen durch. 
    \item Finden Sie eine Diskrepanz? Schieben Sie die Schuld nicht einfach auf die numerische Simulation. Stellen Sie Hypothesen auf, z. B. prüfen Sie 
der Fehler verschwindet mit der Schrittweite.
\end{enumerate}



NUMERICAL PLAUSIBILITY CHECKS
\begin{itemize}
    \item ine numerische Simulation bietet mehr Möglichkeiten zur Überprüfung der Gültigkeit Ihrer Kinematik und Bewegungsgleichungen 
Bewegung. Zum Beispiel:
    \item Calculate kinetic and potential energy of your system, and work done by all forces and 
moments. Check POWER derived from both. This is one of the most powerful checks in finding 
mistakes, so always apply it.
    \item Calculate derivatives (for example velocity of a point, or angular momentum) both analytically 
and numerically. Overlay. 
    \item Check special cases: Derive (on paper) the solution for a much simplified problem, for example 
a 2D version, or a system where diverse parts of the system are massless. Implement a 
separate function with only those simplified equations of motion. Overlay with the complete 
simulation for these specific parameter choices
\end{itemize}

MORE CODING ADVICE
\begin{enumerate}
    \item Verwenden Sie Funktionen anstelle eines langen Skripts. 
    \item Tauschen Sie nur die notwendigen Variablen mit Funktionen aus.
    \item Test each individual function with expected outputs for given inputs. 
    \item Include checks in your functions, for example for expected input variable dimensions. Only 
reduce your code to optimize its performance later.
    \item Comment while you write the code, not later. 
    \item  Always place units in the comments. Do not just place a numerical value for a parameter. 
    \item Create separate functions for your plausibility checks. Include these in the final code package 
(even if they do not need to run all the time, to improve speed).
    \item Make sure you do not over-write pre-defined variables or function names, such as (in Matlab) 
i,j, gamma, etc.
    \item Be aware of “features” of the software that can occlude mistakes. 
  \item  Never rely only on warnings of the software to detect mistakes.
\end{enumerate}


EINIGE SPEZIFISCHE HINWEISE FÜR MATLAB
\begin{itemize}
    \item Verwenden Sie den Befehl "keyboard" und eine "if"-Bedingung, um Ihre Simulation zu einem bestimmten Zeitpunkt t zu beenden, 
zum Beispiel mit ode45. 
\item Befolgen Sie alle Warnungen, die Ihr Code enthält. Überprüfen Sie den orangefarbenen Balken auf der rechten Seite. Beheben Sie alle 
beheben, auch wenn es in Ihren Augen "kein wirkliches Problem" ist
\item Beachten Sie, dass Matlab als "Feature" stehende und liegende Vektoren addiert und eine Matrix erzeugt. Sie erhalten also nicht immer eine Fehlermeldung, wenn die Dimensionen der Vektoren, die Sie addieren, nicht übereinstimmen! 
\item Beachten Sie, dass Matlab es erlaubt, sich auf Elemente in Matrizen mit einem einzigen Index zu beziehen. Sie werden also nicht 
einen Fehler, wenn man A(3) schreibt, wenn A eine 2x3-Matrix ist...
\end{itemize}





%---------------------------------------------------------
\section{LATEX}

\LaTeX ist ein Softwarepaket, das die Benutzung des Textsatzsystems TeX mit Hilfe von Makros vereinfacht. LaTeX liegt derzeit in der Version $2\epsilon$ vor. LaTeX wurde Anfang der 1980er Jahre von Leslie Lamport entwickelt.Der Name bedeutet so viel wie Lamport TeX. Die Entwicklung wurde seit den 1990er Jahren von einer Anzahl Entwicklern weitergeführt. Heute ist LaTeX die populärste Methode, TeX zu verwenden. 
\cite{wikilatex}
%<https://de.wikipedia.org/wiki/LaTeX> 

Im Gegensatz zu anderen Textverarbeitungsprogrammen, die nach dem What-you-see-is-what-you-get-Prinzip funktionieren, arbeitet der Autor mit Textdateien, in denen er innerhalb eines Textes anders zu formatierende Passagen oder Überschriften mit Befehlen textuell auszeichnet. 
\cite{wikilatex}
%Aus <https://de.wikipedia.org/wiki/LaTeX#Grundprinzip> 

Ich benutze Latex zum Erstellen wissenschaftlicher und didaktischer Dokumente (Berichte, Paper, Doktorarbeit, Vorlesungsunterlagen) aber auch zum Erstellen meines Lebenslaufes, Bewerbungsschreiben, formale Briefe im Allgemeinen. Einmal installiert und zum Laufen gebracht, erspart es wirklich sehr viel Formatierungsärger und man kann sich beim Schreiben auf den Inhalt konzentrieren. Die Form passt automatisch.

Um mit LaTeX zu arbeiten, werden zwei Dinge benötigt: Einerseits die LaTeX-Software und andererseits eine Entwicklungsumgebung, mit der der LaTeX-Code eingegeben und die Umsetzung in ein fertig gesetztes Dokument angestoßen wird. Hier mein Paket (alles kostenfrei). Bitte in genau dieser Reihenfolge installieren.

\begin{itemize}
    \item • MiKTeX (TEX- Distribution)
\item Ghostscript (zur PS- Erstellung)
\item GSview (PS- Viewer. Ghostview eignet sich auch zur Betrachtung von pdf Dateien. Es ist in der Lage pdf Dateien zu „reloaden“ wenn sie durch einen LateX Kompiliervorgang erneuert wurden. Acrobat- Reader geht natürlich auch.)
\item JabRef, (Javaprogramm zur Verwaltung von Bibliografie-Dateien (.bib))
\item Inkscape (zum Erstellen von EPS- Grafiken)
\item TeXnicCenter (Editor)
\end{itemize}

Der Wizard legt ein paar Ausgabekonfigurationen an. Er fragt dabei meist auch nach Pfaden zu verschiedenen exe- Dateien. Sie finden sich im MikTeX- Ordner 
%(etwa da: \texttt{:\Program Files\MiKTeX 2.9\miktex\bin\x64 })
Man muss die Ausgabekonfiguration "LaTeX->PS->PDF" auswählen, damit mein Beispielprojekt funktioniert. 

\hyperlink {https://latex.tugraz.at/}{HILFE ONLINE:}
Dort wird das Zitieren auch sehr ausführlich erklärt und man findet Links zu den Downloads der bei der Installation benötigten Software.

OVERFEAF!




Alles zum Thema Latex ist in
\href{https://tobi.oetiker.ch/lshort/lshort.pdf}{T.~Oetiker: The Not So Short Introduction to \LaTeXe}
beschrieben.

%---------------------------------------------------------
\section{Arbeit im Regelungstechniklabor}

\subsection{Rolle der Betreuer\_innen}
Deine Erst- und Zweitprüfer\_innen sind erste Ansprechpartner\_innen bei allen auftretenden Unklarheiten. Neben fachlichen Fragen werden sie dir auch bei der Beantwortung aller anderen Fragen, die im Zusammenhang mit deiner Abschlussarbeit an der HTW auftreten, z. B. Beschaffungen, Arbeitsaufträge an die Werkstätten, etc., helfen,  Die Initiative sollte jedoch stets von dir ausgehen.


\subsection{Arbeitsplatz, Arbeitszeit und Schlüssel im Regelungtechnik- Labor}

Wenn benötigt, lass dir zu Beginn deiner Arbeit vom Laboringenieur einen Arbeitsplatz und, wenn vorhanden, ein abschließbares Schrankfach zuweisen und die Schlüsselregeln  erklären. Bitte akzepiere zu Beginn deiner Arbeit im Labor die aktuellen Hygieneregeln und die Laborordnung mit deiner Unterschrift.

Die Arbeitstische im Regelungstechnik- Labor werden in der Regel von mehreren Personen benutzt. Räumen Sie daher Ihren Tisch nach Benutzung stets auf.

Die Arbeitsräume sind während der regulären Arbeitszeit von 8:00 bis 20:00 Uhr zugänglich. Die Laborordung ist Um dir die Möglichkeit zu geben, in begründeten Ausnahmefällen auch in der Zeit von 20.00 bis 22.00 Uhr noch arbeiten zu können, wird folgende Regelung vereinbart:
\begin{itemize}
    \item In diesem Zeitraum dürfen ohne Betreuer\_in keine Hardware-Arbeiten durchgeführt werden.
    \item Das Haus muss spätestens kurz vor 22.00 Uhr verlassen werden. Ab 22.00 Uhr werden die Haustüre und die Flureingangstüren vom Schließdienst versperrt.
    \item Jeder außerhalb der regulären Arbeitszeit und ohne Betreuer arbeitende Diplomand ist verantwortlich für alle selbst verursachten Schäden. Alle Haftungsansprüche gegenüber dem  Lehrstuhl und den für den Lehrstuhl handelnden Personen sind, soweit gesetzlich zugelassen, ausgeschlossen.
\end{itemize}


\subsection{Werkzeug, Werkstattmaterial, Werkstattaufträge}

Werkzeuge und Werkstattmaterial werden vom Laboringenieur ausgegeben. Werkstattaufträge sind mit entsprechenden Konstruktionszeichnungen/-skizzen mit der Betreuerin/ dem Betreuer abzuklären.

\subsection{Laborgeräte, Gerätebeschreibungen und Bauelemente}

Laborgeräte und Gerätebeschreibungen können Sie beim LÖaboringenieur entleihen. 
%Im Gerätelager muss ein Entnahmeschein  hinterlegt werden, der vom Betreuer zu unterzeichnen ist. 
Alle nicht mehr benutzten Geräte sind sofort in das Lager zurückzugeben. Achten Sie bitte auf eine sachgemäße und pflegliche Behandlung der Geräte. Alle Schäden an den Geräten sind umgehend schriftlich an den Laboringenieur zu melden. Wenn fahrlässiges Verschulden vorliegt, wird die Benutzerin/ der Benutzer zum Schadenersatz herangezogen.

Elektronische Bauelemente werden auch vom Laboringenieur ausgegeben.

\subsection{Einkäufe und Bestellungen}

Einkäufe und Bestellungen sind über den Laboringenieur abzuwickeln. Werden Sie mit Besorgungen betraut, so achten Sie bitte sorgfältig auf Kassenbelege, Lieferscheine und Rechnungen, die alle sofort bei Ihrem Betreuer abzuliefern sind. 

\subsection{Haftung}
Die Haftung des/der Diplomanden/Diplomandin erstreckt sich auf alle Leihgaben, die im abschließbaren Schrankfach untergebracht werden können. 
 %Wir empfehlen  dringend, Werkzeuge, kleinere Messgeräte, Bücher und Bauteile nach Beendigung der Arbeit stets wegzuschließen.



%---------------------------------------------------------%
%TOOLS
%\section{GitHub}
%Versionierungsprogramm
